% Privacy-first section 2. Budget 0.75 page, no figure.
\section{Background and related work}
\label{sec:related}

\paragraph{Memorisation in language models.}
Verbatim extraction from a prefix is the standard operationalisation of memorisation
\citep{carlini2023quantifying}, and planted canaries are the standard way to obtain a ground truth
for it \citep{jagielski2023canary}. Repetition is the best-studied driver
\citep{tirumala2022memorization}, with capacity arguments explaining why
\citep{morris2025capacity}. Two assumptions carry over into genomics without examination. The first
is that a prefix-conditioned decoder measures what a model holds; on masked models it does not,
because the conditional it probes is not the one the model was trained on
\citep{mireshghallah2022mlm,wang2026infilling}. The second is that a model which reproduces nothing
has retained nothing, which \S\ref{sec:ident} contradicts directly, and which
\citet{che2026memorization} show cannot hold in general: counterfactual memorisation and adaptive
extraction do not bound one another in either direction.

\paragraph{Tokenisation as a cause.}
That byte-pair vocabularies increase memorisation in text was shown by
\citet{kharitonov2021bpe}, who attribute the effect to the shorter token sequences a larger
vocabulary produces. In text the two cannot be separated: changing the vocabulary changes what a
token \emph{means}, and the vocabulary co-adapts to the corpus. DNA is the setting where the axis is
clean --- a closed four-letter alphabet means every tokenizer is a lossless view of identical
content, and sequence length can be held fixed while granularity varies. Our results support the
dependence on token granularity and locate the mechanism elsewhere than
\citet{kharitonov2021bpe} propose: not in sequence length, but in how many positions must be decoded
correctly at once (\S\ref{sec:slots}).

\paragraph{Genomic privacy.}
Deciding whether a given person is in a cohort, from an artefact the cohort released, is the oldest
attack in this literature. It has been demonstrated on allele-frequency releases
\citep{homer2008resolving}, on beacon services that answer only yes or no
\citep{shringarpure2015beacon,vonthenen2019reidentification}, on polygenic scores
\citep{paige2021reconstructing,nikitin2026prs}, and on summary statistics in general
\citep{li2013membership,tramer2015bounded}. Two things carry over to us. The harm of record in this
field is identification rather than reproduction \citep{erlich2014routes}, which is the opposite
emphasis to the memorisation literature above. And the adversary here sees weights, a newer object
and a stronger one: queryable without limit, and fitted to the individuals rather than to a summary
of them. Generative models trained on genomes have been audited for membership leakage in their
\emph{output} \citep{oprisanu2022synthetic,belfiore2025bihmia}; we audit released weights.

\paragraph{The study this extends.}
\citet{nemecek2026quantifying} fine-tunes four released genomic models --- DNABERT-2, HyenaDNA, a
12.9M \texttt{SimpleDNALM} and Evo --- on four corpora, and reports that they memorise to very
different degrees. It does not isolate a cause, and its comparison varies architecture, parameter
count, pretraining objective and tokenizer together. We hold the first three fixed where we can,
use its \texttt{SimpleDNALM} as our controlled backbone so the two sets of numbers are comparable,
and add the measurement machinery that \S\ref{sec:measure} argues the question requires. Where that
study reports a single extraction rate per model, we report what each of three decoders reads from
the same weights, which is what revealed that the ordering depends on the instrument.

\paragraph{Scoring.}
Pseudo-log-likelihood \citep{salazar2020pll} gives a masked model a sequence score an adversary can
compute without the answer, which is what lets a search over reconstructions choose among its own
fixed points. Recent work on infilling and diffusion decoders
\citep{wang2026infilling,luo2026diffusion} and on the dynamics of memorisation during training
\citep{vonarx2026dynamics} is concurrent with ours and complementary: it asks how to decode better,
where we ask what a decoder's failure licenses one to conclude.
